
\documentclass[10pt]{article}
\usepackage[utf8]{inputenc}
\usepackage[T1]{fontenc}
\usepackage{amsmath, amssymb, xcolor, geometry, enumitem, multicol, tcolorbox}

\definecolor{mainblue}{RGB}{0, 102, 204}
\geometry{a4paper, margin=1.2cm, top=1cm, bottom=3cm, footskip=1.2cm}

\begin{document}
\enlargethispage{2.5cm}

% Modern Header
\begin{tcolorbox}[colback=mainblue!5, colframe=mainblue, sharp corners, boxrule=0.5pt]
    \small \textbf{Unit:} undefined \hfill \textbf{Name:} \rule{3cm}{0.4pt} \\
    \small \textbf{Topic:} Variables and Expressions / Order of Operations \hfill \textbf{Date:} \rule{2cm}{0.4pt} \\
    \centering \textbf{\large Foundation (Level -1)}
\end{tcolorbox}

\vspace{0.2cm}

% MCQ Section
\begin{multicols}{2}
\begin{enumerate}[label=\textbf{Q\arabic*.}, leftmargin=*, itemsep=6pt, parsep=0pt]

  \item \small What is the variable in the expression $2x$?
  \begin{enumerate}[label=(\Alph*), leftmargin=0.6cm, nosep, itemsep=2pt]
    \item 2
    \item $x$
    \item 2$x$
    \item None
    \item Expression
  \end{enumerate}

  \item \small Identify the constant in the expression $4 + 3x$
  \begin{enumerate}[label=(\Alph*), leftmargin=0.6cm, nosep, itemsep=2pt]
    \item 3
    \item 4
    \item $3x$
    \item None
    \item Expression
  \end{enumerate}

  \item \small What is the coefficient of $x$ in the expression $5x$?
  \begin{enumerate}[label=(\Alph*), leftmargin=0.6cm, nosep, itemsep=2pt]
    \item 1
    \item 5
    \item $x$
    \item None
    \item Expression
  \end{enumerate}

  \item \small Which of the following is an example of a variable: $x$, $2$, or $3y$?
  \begin{enumerate}[label=(\Alph*), leftmargin=0.6cm, nosep, itemsep=2pt]
    \item 2
    \item $x$
    \item $3y$
    \item $x$ or $3y$
    \item None
  \end{enumerate}

  \item \small Identify the operator in the expression $7 - 2$
  \begin{enumerate}[label=(\Alph*), leftmargin=0.6cm, nosep, itemsep=2pt]
    \item 7
    \item 2
    \item -
    \item None
    \item Expression
  \end{enumerate}

  \item \small What does the term $2x$ represent in the expression $2x + 5$?
  \begin{enumerate}[label=(\Alph*), leftmargin=0.6cm, nosep, itemsep=2pt]
    \item Constant
    \item Variable
    \item Coefficient
    \item Expression
    \item None
  \end{enumerate}

  \item \small Is $3y$ a variable or an expression?
  \begin{enumerate}[label=(\Alph*), leftmargin=0.6cm, nosep, itemsep=2pt]
    \item Variable
    \item Expression
    \item Constant
    \item Coefficient
    \item None
  \end{enumerate}

  \item \small What is the term for a number that is multiplied by a variable, such as $4$ in $4x$?
  \begin{enumerate}[label=(\Alph*), leftmargin=0.6cm, nosep, itemsep=2pt]
    \item Constant
    \item Coefficient
    \item Variable
    \item Expression
    \item None
  \end{enumerate}

\end{enumerate}
\end{multicols}

\vspace{-0.2cm}
\hrule
\vspace{0.2cm}
\textbf{\small Open Ended Questions (FRQ)}
\begin{enumerate}[label=\textbf{Q\arabic*.}, start=9, itemsep=5pt, topsep=0pt]

  \item \small Draw a diagram to illustrate the difference between a constant and a variable in an algebraic expression. Provide a brief explanation. \vspace{2cm}

  \item \small Explain, using examples, how the order of operations is used to simplify expressions. Be sure to include at least two expressions and the step-by-step simplification process for each. \vspace{4cm}

\end{enumerate}

\vfill

% Chef's Tips Footer
\begin{tcolorbox}[colback=blue!5, colframe=mainblue!50, title=\small \textbf{Chef's Tips}, fonttitle=\bfseries, bottom=1mm]
\begin{itemize}[noitemsep, topsep=2pt, leftmargin=0.5cm]
  \item \scriptsize Emphasize the distinction between variables, constants, and expressions to help students understand the fundamental concepts of algebra.\item \scriptsize Use visual aids, such as diagrams and graphs, to illustrate the relationships between variables, constants, and expressions.\item \scriptsize Provide students with numerous examples and opportunities for practice to reinforce their understanding of the order of operations and how to simplify expressions.
\end{itemize}
\end{tcolorbox}

\end{document}
  